\defmodule {LATBASIS}

This module offers tools to manipulate lattice bases
(see Section 1.3).
Each lattice is represented by a basis $V$ and its dual $W$.
It is sometimes possible, as in the case with lattices associated with
LCGs or MRGs, to multiply a lattice (and its dual) by a constant factor
in such a way that they are included in $\Z^t$, allowing exact representation
of basis vector coordinates.
The duality relation will now read $V_i\cdot W_j = m\delta_{ij}$
for some integer constant {$m$}.

\iffalse  %%%%%%%%%%%%%
There is a choice between two representations for the vector coordinates:
{\tt LONGREAL} or {\tt SuperInteger}.
The latter is a type defined in a package named SENTIERS which
offers tools for manipulating arbitrary large integers.
The {\tt SuperInteger} representation should be used
when the coordinates of the (dual) basis vectors, are too large
to fit a {\tt LONGREAL} without loss of precision.
But computations are then much slower.
The program {\tt select} decides which one is used.
All computations will be performed using appropriate procedures,
depending on the representation.
\fi  %%%%%%%%%%%%%%%%%

The squared lengths of the basis vectors and that of their duals
are stored as {\tt LONGREAL} in {\tt VVLR} and {\tt WWLR} respectively.
However, if they are too large to fit in a {\tt LONGREAL},
even approximately, then {\tt SISquares} is set to {\tt TRUE} and
they are stored as {\tt SuperInteger} in {\tt VVSI} and {\tt WWSI}.  

A variable of type {\tt Basis} is created by calling {\tt CreateBasis}.
The squared vector lengths {\tt VVLR} and {\tt WWLR} 
(or {\tt VVSI} and {\tt WWSI})
can then be updated by calling {\tt UpdateVVWW}.
But it is not necessary to update them before calling any of the procedures
of the modules {\tt REDBAS} and {\tt REDBAS2}
(including {\tt ReductMinkowski} and {\tt ShortestVector}).

\bigskip\hrule
\code

DEFINITION MODULE LATBASIS;

FROM SUPINT  IMPORT SuperInteger;

IMPORT BASISLR, BASISSI;
FROM BASIS    IMPORT BMatr, BScal;
FROM CONFIG   IMPORT BDim;
FROM NORM     IMPORT NormType;
FROM MULT     IMPORT MScal;

TYPE
   Basis = POINTER TO InfoBasis;
   InfoBasis = RECORD
      ResDim     : BDim;             (* Dimension of allocated memory.      *)
      Dim        : BDim;             (* Actual dimension.                   *)
      mLR        : LONGREAL;         (* Modulus m.                          *)
      mSI        : SuperInteger;     (* Modulus m.                          *)
      mLR2       : LONGREAL;         (* m^2, rescaled if SISquares.         *)
      mSI2       : SuperInteger;     (* Value of m^2 when SISquares.        *)
      V, W       : BMatr;            (* Basis and dual basis vectors.       *)
      Norm       : NormType;
      Lmin, Lmin2  : LONGREAL;       (* Length and squared length of shor-  *)
                                     (* test basis vector, using Norm.      *)
                                     (* Different meaning in REDBB.         *)
      SISquares    : BOOLEAN;
      VVLR, WWLR   : BASISLR.BVect;  (* Squared Euclidean vector lengths,   *)
                                     (* if not SISquares.                   *)
      EchVV, EchWW : SuperInteger;   (* Scaling factors when SISquares.     *)
      VVSI, WWSI : BASISSI.BVect;    (* Rescaled squared Euclidean vector   *)
                                     (* lengths, if SISquares.              *)
      XX : ARRAY BDim OF BOOLEAN;    (* Used for ReductMinkowski only.      *)
                                     (* Ask R. Couture for details.         *)
      END;
\endcode
  \tabb
   A lattice basis.  
   {\tt ResDim} is the dimension of the arrays allocated for the basis
   (the maximum value that {\tt Dim} can take).
   The memory for the {\tt BVect} and {\tt BMatr} elements is allocated
   for dimensions 0 to {\tt ResDim} if {\tt ResDim} $> 0$, whereas
   no memory is allocated if {\tt ResDim} $= 0$.

   {\tt Dim} is the current dimension of the basis,
   {\tt V} contains the basis, and {\tt W} the $m$-dual basis. 
   The integer constant $m$ used in the duality relation is stored
   as a {\tt LONGREAL} in {\tt mLR}, and as a {\tt SuperInteger}
   in {\tt mSI}.

   {\tt Norm} is the type of norm used to measure the vector lengths.
   {\tt Lmin} and {\tt Lmin2} are used to store the length and the 
   squared length of the shortest vector found to date, using this norm.
   (Usually, only one of these 2 variables is used at a time; 
   see {\tt REDBAS}.)

   The boolean {\tt SISquares} is set to {\tt TRUE} by {\tt InitBasis}
   when $m$ exceeds approximately $2^{480}$, and to {\tt FALSE} otherwise.

   When {\tt SISquares = FALSE}, {\tt mLR2} contains the value of $m^2$,
   {\tt VVLR} and {\tt WWLR} {\em may\/} memorize the squared Euclidean 
   lengths of the vectors in $V$ and in $W$. 
   They are not always updated, and when they are not, they
   are set to negative values.

   When {\tt SISquares = TRUE}, {\tt mSI2} contains $m^2$,
   {\tt mLR2} contains $m^2$ divided (rescaled) by {\tt EchVV * EchWW},
   the squared Euclidean vector lengths (again, not always updated),
   are stored in {\tt VVSI} and {\tt WWSI}, and 
   {\tt Lmin2} is rescaled by {\tt EchVV}.
   Also, the squared vector lengths rescaled by the factors
   {\tt EchVV} and {\tt EchWW}, are put in {\tt VVLR} and {\tt WWLR},
   respectively, by {\tt SetScale}.
   The scaling factors {\tt EchVV} and {\tt EchWW} are computed
   by the procedure {\tt SetScale} and set to values such that
   the log of the squared Euclidean length of the longest vector is 
   not too far from 0.  
   The aim of this rescaling is to avoid {\tt LONGREAL} overflow.
   \endtabb
%\hide
%VAR
%   (* This has been moved to  BASIS^.XX  in LATBASIS.  *)
%   XX  : ARRAY BDim OF BOOLEAN;    (* ??? *)    
%\endhide
\code

PROCEDURE CreateBasis (VAR B : Basis);
\endcode
 \tab Creates the data structure {\tt B} for a basis. 
    {\tt SISquares} is initialized to {\tt FALSE},
    {\tt Norm\/} is initialized to {\tt L2Norm},
    and {\tt ResDim} is set to 0.
    One must then call {\tt InitBasis} to reserve the memory for
    the {\tt BVect} and {\tt BMatr} elements.
 \endtab
\code

PROCEDURE InitBasis (B : Basis;  m : MScal;  Nor : NormType;  d : BDim);
\endcode
   \tab
    Initializes the modulus and the norm for the basis {\tt B}.  
    Reserves the memory for the {\tt BVect} and {\tt BMatr}
    fields of {\tt B}, for indices 0 to $d$ if $d>0$.
    Here, {\tt d} is the new value of {\tt ResDim}.
    If {\tt d < B\^{}.ResDim}, then the memory already reserved which
    is no longer needed is deallocated.
    \endtab
\code

PROCEDURE DeleteBasis (B : Basis);
\endcode
   \tab  Does the same as {\tt InitBasis (B, B\^{}.m, B\^{}.Norm, 0)}, 
    and then the  opposite of {\tt CreateBasis (B)}.
%   Deallocates the memory allocated for {\tt V}, {\tt W}, {\tt VVSI} 
%   and {\tt WWSI}, at a previous call to {\tt InitBasis}. 
%   Must be called only if {\tt InitBasis} has been called previously.
   \endtab
\code

PROCEDURE WriteBasis (B : Basis);
\endcode
  \tab  Writes content of {\tt B} in current output file.
  \endtab
\code

PROCEDURE ReadBasis (VAR B : Basis);
\endcode
  \tab  Reads a basis from current input file.
        The basis must have been written to the file by {\tt WriteBasis}.
  \endtab
\code

PROCEDURE CopyBasis (B1, B2 : Basis);
\endcode
 \tab  Copies {\tt B1} into {\tt B2}.  \endtab
\code

PROCEDURE Permute (B : Basis;  i, j : BDim);
\endcode
   \tab Exchanges vectors i and j in the basis and its dual. \endtab
\code

PROCEDURE SortBasis (B : Basis;  d : BDim);
\endcode
   \tab  Sorts the basis vectors (with index from $d+1$ to {\tt B\^{}.Dim}) 
   by increasing length.
   The dual vectors are permuted accordingly.
   Assumes that {\tt VVLR} and {\tt WWLR} 
   (or {\tt VVSI} and {\tt WWSI}, if {\tt SISquares}) are up to date.
%  Note: no longer updates {\tt VVLR} (or {\tt VVSI}). 
 \endtab
\code

PROCEDURE Dualize (B : Basis);
\endcode
  \tab Exchanges {\tt V, VVLR, VVSI, EchVV} and {\tt W, WWLR, WWSI,
       EchWW}.
 \endtab
\code

PROCEDURE UpdateVVWW (B : Basis;  d : BDim);
\endcode
  \tab  Updates {\tt VVLR} and {\tt WWLR} (or {\tt VVSI} and {\tt WWSI}
     if {\tt SISquares}) as the square Euclidean lengths of the vectors
     in {\tt V} and {\tt W} (for vectors 
     with index from $d+1$ to {\tt B\^{}.Dim}).
 \endtab
\code

PROCEDURE EraseVVWW (B : Basis);
\endcode
  \tab
    Sets the vectors {\tt VVLR} and {\tt WWLR} (or {\tt VVSI} and {\tt WWSI}
    if {\tt SISquares}) to negative values. 
\endtab
\code

PROCEDURE CheckDuality (B : Basis): BOOLEAN;
\endcode
  \tab
   Check that {\tt Basis B} satisfies the duality relation
   $V[i]\cdot W[j] = m\,\delta_{ij}$. If so, returns {\tt TRUE}, otherwise 
   returns {\tt FALSE}.
\endtab
\code

PROCEDURE BaseEquivalence (B1, B2 : Basis) : BOOLEAN;
\endcode
  \tab
   Check that {\tt Basis B1} and {\tt B2} are equivalent.
   If so, returns {\tt TRUE}, otherwise returns {\tt FALSE}.
\endtab
\code

PROCEDURE RandomLattice (B : Basis;   m : BScal;   d : BDim);
\endcode
 \tab Constructs a {\em random\/} basis for an integral lattice,
   chosen with the uniform distribution over the {\em finite\/} set
   of all lattices whose fundamental volume is $m^{d-1}$ and which
   contains $m \Z^d$ as a sublattice, for the specified values of $m$ and $d$.
 \endtab
\code

END LATBASIS.
\endcode
